\begin{helper}
Fix \(\eta\) and \(n\), and write
\[
L=\log n,\qquad
k_R=(1+\eta)\sqrt{nL},\qquad
K=\noderef{TwoBiteNaturalIndependenceScale}(1+\eta,n),\qquad
k=(K:\mathbb R).
\]
Assume
\[
1<L,\qquad k_R\le k,\qquad 16L^2\sqrt L\le k_R.
\]
Then
\[
\frac1{2L}\noderef{RealChooseTwo}(k)+2kL^2+
\frac1{\sqrt L}\noderef{RealChooseTwo}(k)
\le
\frac2{\sqrt L}\noderef{RealChooseTwo}(k).
\]
\end{helper}
\begin{proof}
Since \(K\) is a natural number, \(k\ge0\).  From \(1<L\) we have
\(\sqrt L\ge1\) and \(L^2\ge1\), so
\[
2\le16L^2\sqrt L\le k_R\le k.
\]
Applying \(\noderef{RealChooseTwoQuadraticBounds}\) to this \(k\) gives
\[
\frac{k^2}{4}\le\noderef{RealChooseTwo}(k),
\]
and in particular \(\noderef{RealChooseTwo}(k)\ge0\).

Because \(1<L\), also \(\sqrt L\le L\), hence
\[
\frac1{2L}\noderef{RealChooseTwo}(k)
\le
\frac1{2\sqrt L}\noderef{RealChooseTwo}(k).
\]
The threshold \(16L^2\sqrt L\le k_R\le k\), multiplied by \(k\ge0\), gives
\[
16kL^2\sqrt L\le k^2.
\]
Dividing by \(8\sqrt L>0\) gives
\[
2kL^2\le\frac{k^2}{8\sqrt L}
=\frac1{2\sqrt L}\frac{k^2}{4}
\le\frac1{2\sqrt L}\noderef{RealChooseTwo}(k).
\]
Adding these two estimates and the remaining
\(\frac1{\sqrt L}\noderef{RealChooseTwo}(k)\) term yields exactly the claimed
three-term comparison.
\end{proof}
